%! AP Statistics — Unit 3, Lesson 3.2 — Activity KEY
\documentclass[10pt]{article}
\usepackage{apstats-article}
\usepackage{apstats-key}

\begin{document}

\pageheader{Unit 3, Lesson 3.2}{Group Activity KEY: Can We Conclude That?}
\namepartnerperiod

\vspace{0.1in}

\textbf{Study 1:} TV watching and BMI.
\begin{practicebox}
\small
\begin{enumerate}[leftmargin=1.5em, itemsep=6pt]
  \item \ans{Observational Study}
  \item Population: \ans{All adults.} \quad Sample: \ans{The 1,000 randomly selected adults.}
  \item \ans{Yes — the adults were randomly selected, so results can generalize to the adult population.}
  \item \ans{No — this is an observational study; we cannot establish causation. Researchers did not impose TV-watching.}
  \item \ans{Possible confounders: sedentary lifestyle, diet, socioeconomic status (any one acceptable).}
\end{enumerate}
\end{practicebox}

\vspace{0.08in}

\textbf{Study 2:} Reading program experiment.
\begin{practicebox}
\small
\begin{enumerate}[leftmargin=1.5em, itemsep=6pt]
  \item \ans{Experiment}
  \item Random assignment: \ans{Yes.} \quad Random selection: \ans{No (volunteers at one school).}
  \item \ans{Yes — random assignment was used, balancing confounding variables. The program caused the score difference.}
  \item \ans{No — students were volunteers at one school, not randomly selected. We cannot generalize to all students.}
\end{enumerate}
\end{practicebox}

\vspace{0.08in}

\textbf{Study 3:} Diet and cholesterol observational study.
\begin{practicebox}
\small
\begin{enumerate}[leftmargin=1.5em, itemsep=6pt]
  \item \ans{Observational Study}
  \item \ans{Yes.}
  \item Generalize to county: \ans{Yes.} \quad To all U.S. adults: \ans{No} (only sampled from one county).
  \item \ans{No — observational study; cannot establish causation.} Justify: \ans{Confounding variables such as exercise, age, or genetics may explain the association.}
\end{enumerate}
\end{practicebox}

\vspace{0.08in}

\begin{practicebox}
\small
\textbf{Group Synthesis:}

\renewcommand{\arraystretch}{1.6}
\begin{tabularx}{\linewidth}{@{} >{\bfseries}p{0.38\linewidth} X X @{}}
 & \textbf{Causation?} & \textbf{Generalize?} \\
\hline
Obs. study, random sample & \ans{No} & \ans{Yes (to population)} \\
Experiment, rand. assign., no rand. sample & \ans{Yes} & \ans{No (sample only)} \\
Experiment, rand. assign., rand. sample & \ans{Yes} & \ans{Yes (to population)} \\
\end{tabularx}
\end{practicebox}

\end{document}
